%
% Qualificacao_Doutorado.tex (LateX)
% 
% Objetivo: Arquivo principal do relatório de qualificação de doutorado.
% baseado em um template para a geração de documentos em LaTeX.
% 
% Versão 1.0
% 
% Site: http://www.dirackslounge.online
% 
% Programador: 
%		(1.4) - Lauro César Araujo 
%			Template, distribuição e manutenção
%		(1.5) - Rodolfo A. C. Neves (Dirack) 07/10/2019 
%			Modificações e utilização neste relatório
% 
% Email: rodolfo_profissional@hotmail.com
% 
% Licença (Versão modificada): GPL-3.0 <https://www.gnu.org/licenses/gpl-3.0.txt>.
%
% Licença (Versão original): LaTeX Project Public License (LPPL - 1.3) <http://www.latex-project.org/lppl.txt>
%
% Documentação extra: <http://abntex2.googlecode.com/>

\documentclass[
	% -- opções da classe memoir --
	12pt,				% tamanho da fonte
	openright,			% capítulos começam em pág ímpar (insere página vazia caso preciso)
	oneside,			% para impressão em verso e anverso. Oposto a oneside
	a4paper,			% tamanho do papel. 
	% -- opções da classe abntex2 --
	%chapter=TITLE,		% títulos de capítulos convertidos em letras maiúsculas
	%section=TITLE,		% títulos de seções convertidos em letras maiúsculas
	%subsection=TITLE,	% títulos de subseções convertidos em letras maiúsculas
	%subsubsection=TITLE,% títulos de subsubseções convertidos em letras maiúsculas
	% -- opções do pacote babel --
	english,			% idioma adicional para hifenização
%	french,				% idioma adicional para hifenização
%	spanish,			% idioma adicional para hifenização
	brazil				% o último idioma é o principal do documento
	]{abntex2}

\usepackage{multirow}
\usepackage{amsmath}
\usepackage{tocloft}
\usepackage{cmap}			% Mapear caracteres especiais no PDF
\usepackage{lmodern}			% Usa a fonte Latin Modern			
\usepackage[T1]{fontenc}		% Seleção de códigos de fonte.
\usepackage[utf8]{inputenc}		% Determina a codificação utiizada (conversão automática dos acentos)
\usepackage{makeidx}            	% Cria o indice
\usepackage{amssymb,amsfonts,amsmath,wasysym}
\usepackage{lastpage}			% Usado pela Ficha catalográfica
\usepackage{indentfirst}		% Indenta o primeiro parágrafo de cada seção.
\usepackage{nomencl} 			% Lista de simbolos
\usepackage{color}			% Controle das cores
\usepackage{graphicx}			% Inclusão de gráficos
\usepackage{microtype} 			% para melhorias de justificação
\usepackage{subfig}
\usepackage{float} 			% Colocar a figura no local certo
\usepackage{scalefnt} 			% Pacote da redimensionar a fonte de tabelas, figuras e equações.
\usepackage{placeins}
\allowdisplaybreaks
\usepackage{kantlipsum}
\usepackage{pdfpages}
\usepackage{titlesec}
\usepackage[brazilian,hyperpageref]{backref}	 % Paginas com as citações na bibl
\usepackage[alf]{abntex2cite}			 % Citações padrão ABNT
\usepackage{tensor}
\usepackage{setspace}
\usepackage{caption}
\usepackage{amsmath}
\usepackage{enumerate}
\usepackage[portuguese,ruled,lined]{algorithm2e}
\usepackage{algorithmic}

\renewcommand{\figurename}{Figura-}
\usepackage[figurename=Figura]{caption}

\renewcommand{\chapnumfont}{\bfseries}
\renewcommand{\ABNTEXfontereduzida}{\mdseries\footnotesize}
\renewcommand{\ABNTEXchapterfontsize}{\bfseries \normalsize}
\renewcommand{\ABNTEXsectionfontsize}{\normalsize}
\addto{\captionsbrazil}{\renewcommand{\contentsname}{\textbf{SUMÁRIO}}}
\addto{\captionsbrazil}{\renewcommand{\bibname}{\bfseries \textbf{REFERÊNCIAS}\selectfont}}
\renewcommand{\apendicesname}{\bfseries \selectfont  APÊNDICES}
\renewcommand{\apendicename}{\bfseries\selectfont APÊNDICE}
\renewcommand{\anexosname}{\bfseries ANEXOS}

\renewcommand*{\backrefalt}[4]{
	\ifcase #1
	
	\or
	
	\else

	\fi}

%% Informações básicas da CAPA
\instituicao{
  Universidade Federal do Pará -- UFPA
  \par
  Instituto de Geociências
  \par
  Programa de Pós-Graduação em Geofísica}
\titulo{INVERSÃO DO MODELO DE VELOCIDADES NO DOMÍNIO ERC UTILIZANDO AS APROXIMAÇÕES DE TEMPO DE TRÂNSITO SRC NÃO HIPERBÓLICO}

\autor{Rodolfo André Cardoso Neves}
\local{Belém-Pará}
\data{2021}

\orientador{Prof. Dr. João Carlos Ribeiro Cruz}
\coorientador{}

\preambulo{Relatório apresentado ao Programa de Pós-Graduação em Geofísica do Instituto de Geociências 
da Universidade Federal
do Pará, em cumprimento às exigências para obtenção do grau de Doutor em Geofísica.}

\definecolor{blue}{RGB}{41,5,195}
\definecolor{black2}{RGB}{39,64,139}

%% Configurações padrão do PDF
\hypersetup{
		backref=true,
		pagebackref=true,
		bookmarks=true,         		% show bookmarks bar?
		pdftitle={\imprimirtitulo}, 
		pdfauthor={\imprimirautor},
    	pdfsubject={\imprimirpreambulo},
		pdfkeywords={PALAVRAS}{CHAVES}{abnt}{abntex}{abntex2},
	    pdfproducer={LaTeX with abnTeX2}, 		% producer of the document
	    pdfcreator={\imprimirautor},
    	colorlinks=true,       				% false: boxed links; true: colored links
    	linkcolor=black,          			% color of internal links
    	citecolor=black,        			% color of links to bibliography
    	linkcolor=black,          			% color of internal links
    	citecolor=black,        			% color of links to bibliography
    	filecolor=black,      				% color of file links
		urlcolor=black,
		bookmarksdepth=4
}

% Indentação do parágrafo
\setlength{\parindent}{1.3cm}

% Espaçamento entre parágrafos
\setlength{\parskip}{0.2cm}

% Espaçamento entre linhas
\OnehalfSpacing	

% compilar indice
\makeindex

% Compilar lista de abreviaturas e siglas
\makenomenclature

\renewcommand{\sin}{\mathrm{sen}}
\newcommand{\disp}{\displaystyle}
\newcommand{\mbf}{\mathbf}

\hyphenation{geo-fí-si-co}
\hyphenation{MCSEM}

\captionsetup{labelsep=period,font=small,justification=justified,labelfont=md,format=plain,labelsep=endash}

\renewcommand{\imprimircapa}{
\begin{capa}
	\begin{figure}[!t]
		\begin{center}
			\includegraphics[scale=0.3]{images/Logo_UFPA2.pdf}
		\end{center} 
	\end{figure}
			\begin{center}
				{\ABNTEXchapterfont\bfseries{UNIVERSIDADE FEDERAL DO PARÁ \\ INSTITUTO DE GEOCIÊNCIAS \\ PROGRAMA DE PÓS-GRADUAÇÃO EM GEOFÍSICA} }
		\end{center}

\center

\vspace*{1cm}
{\ABNTEXchapterfont\bfseries\large\MakeUppercase{\imprimirautor}}\\
\vspace*{2cm}
{\ABNTEXchapterfont\bfseries\large\imprimirtitulo}\\
\vspace*{2cm}{\ABNTEXchapterfont\bfseries\large{RELATÓRIO DO PEDIDO DE PRORROGAÇÃO DE TESE DE DOUTORADO}}
\vspace*{\fill}\\
{\large\MakeUppercase{\imprimirlocal}}
\par
{\large\imprimirdata}
\vspace*{1cm}
\end{capa}
}


\makeatletter
\renewcommand{\folhaderostocontent}{
\begin{center}
	\vspace*{1cm}
	{\ABNTEXchapterfont\bfseries\large\MakeUppercase{\imprimirautor}}\\
	\vspace*{\fill}\vspace*{\fill}
	{\ABNTEXchapterfont\bfseries\large\imprimirtitulo}\\
	 \vspace{\baselineskip}
	{\ABNTEXchapterfont\bfseries\large{RELATÓRIO DO PEDIDO DE PRORROGAÇÃO DE DOUTORADO}}
	\vspace{\baselineskip}
	\vspace{\baselineskip}
	\vspace*{\fill}
	\abntex@ifnotempty{\imprimirpreambulo}{
	\hspace{5cm}
	\begin{minipage}{10cm}
	\SingleSpacing
	\imprimirpreambulo\\ \\
	{\imprimirorientadorRotulo~Prof. Dr. João Carlos Ribeiro Cruz}\\
	\end{minipage}
	\vspace*{\fill}
	}
	\abntex@ifnotempty{\imprimircoorientador}{
	{\large\imprimircoorientadorRotulo~\imprimircoorientador}
	}
	\vspace*{\fill}
	{\large\imprimirlocal}
	\par
	{\large\imprimirdata}
	\vspace*{1cm}
\end{center}

}

\makeatother

\begin{document}

\imprimircapa

\imprimirfolhaderosto*

\begin{dedicatoria}
   \vspace*{\fill}
   \vspace*{13cm}
   \centering
   \noindent
   \begin{flushright} Dedico este trabalho \`a minha fam\'ilia.
   \end{flushright}
     \vspace*{\fill}
\end{dedicatoria}

\begin{agradecimentos}[\fontsize{12pt}{\baselineskip}\textbf{AGRADECIMENTOS}]
\vspace*{1.5cm}

Ao Prof. Dr. João Carlos pela orientação desta tese, e pela proposta do tema. Além
do suporte e paciência em responder as minhas dúvidas, e das sugestões inteligentes na solução de problemas
que foram aparecendo no caminho.

Agradeço ao professor Sergey Fomel, que apesar de não conhecer pessoalmente, produziu trabalhos
que inspiraram o tema, e disponibilizou gratuitamente
vários
dos programas aqui utilizados.

Agradeço à minha mãe Regina de Nazaré,
à minha irmã Rebeca Cristina, às minhas sobrinhas Ágatha e Anadora, 
e ao meu pai Ricardo Neves, por todo apoio e dedicação
durante a árdua caminhada para a realização deste sonho!

E agradeço à amizade de amigos que conquistei durante o curso de Geofísica, e que de forma direta ou indireta
me ajudaram na realização do trabalho:
Leonardo Reis, Hugo Souza, Diogo Rezende, Antônio Rizimar e Raphael Di Carlo.

\end{agradecimentos}

\begin{epigrafe}
    \vspace*{\fill}
	\begin{flushright}
		\textit{``Todas as coisas excelentes são tão difíceis quanto raras''. \\
		(Baruch Spinoza)}
	\end{flushright}
\end{epigrafe}

\begin{resumo}[\fontsize{12pt}{\baselineskip}\textbf{RESUMO}]
\OnehalfSpacing

Neste trabalho, propomos uma metodologia de inversão do modelo de velocidades pós empilhamento, a partir dos
parâmetros do SRC de afastamento nulo $R_N$, $R_{NIP}$ e $\beta_0$ e da seção empilhada, obtidos no empilhamento ERC.
O método do elemento de reflexão comum (ERC) possui a vantagem de 
não necessitar do conhecimento a priori do modelo de velocidades na obtenção da seção empilhada.
Apenas a velocidade próxima à superfície é requerida. 
Este método proporciona parâmetros importantes para a obtenção do
macromodelo de velocidades na etapa de inversão subsequente.
Utilizamos a aproximação do SRC não hiperbólico para
obter os parâmetros do SRC de afastamento nulo através da otimização global com o algoritmo Very Fast Simulated Aneeling (VFSA),
e calculamos a curva de empilhamento aplicando a condição SDC ($R_N=R_{NIP}$) à esta aproximação.
Obtemos as famílias ERC através da regularização dos dados, aumentando a amostragem das seções de afastamento constante
com a interpolação com filtros adaptativos de predição de erro (FPE) e calculamos as trajetórias ERC 
no plano Ponto médio comum x afastamento em função de dos parâmetros $R_{NIP}$ e $\beta_0$.
A inversão do modelo de velocidades é realizada em um laço iterativo e a atualização do modelo é feita com o VFSA.
O modelo em cada iteração é obtido a partir de respostas de difração simuladas na seção empilhada ERC. Utilizando uma
medida de focalização e a continuação de velocidades, obtemos o modelo de velocidades que melhor focaliza as hipérboles
de difração simuladas. O critério de convergência do modelo de velocidades é a coerência entre os dados modelados
com o modelo de velocidades da iteração atual e os dados pré empilhados.
\vspace{\onelineskip} 
 \noindent
 \par Palavras-chave: Inversão do modelo de velocidades. Aproximações não hiperbólicas do tempo de trânsito SRC. 
 Empilhamento Elemento de Reflexão Comum (ERC). 
\end{resumo}

\begin{resumo}[\fontsize{12pt}{\baselineskip}\textbf{ABSTRACT}]
\OnehalfSpacing

In this work, we propose a new methodology of post stack velocity model inversion, from zero offset CRS parameters
$R_N$, $R_{NIP}$ e $\beta_0$ and stacking section obtained during CRE stacking process.
The main advantage of the CRE method is it does not need a priori knowledge of the macro velocity model to obtain the
stacking section. Only near surface velocity information is required.
This method brings important parameters to the macrovelocity model in the inversion step.
We use the non hyperbolic CRS approximation to obtain the zero offset 
CRS parameters through the Very Fast Simulated Aneeling (VFSA)
global optimization algorithm, and calculate the stacking curve applying the 
CDS condition ($R_N=R_{NIP}$) in this approximantion.
We obtain the CRE Gathers through data regularization, increasing the constant offset gathers sampling with
non adaptative prediction error filters interpolation (PEF) and calculating the seismic traces in the CRE trajectory in the 
CMP x Offset plane in function of the $R_{NIP}$ and $\beta_0$ parameters.
The velocity model inversion is done inside an iterative loop and the model updating uses VFSA algorithm.
In each iteration the velocity model is obtained through diffraction response simulation in the CRE stacked section.
Using the local varimax measure as a focusing criteria and velocity continuation, we obtain the best velocity model
that focus the simulated diffraction hyperbole.
The convergence criteria of the velocity model is the semblance between the modeled data with the 
actual iteration velocity model and the pre-stack seismic data.
\vspace{\onelineskip} 
\noindent 
 \par Palavras-chave: Velocity model inversion. Non-hyperbolic CRS traveltime approximation. 
 Common reflecting element (CRE) stacking. 
\end{resumo}

\renewcommand{\listfigurename}{\fontsize{12pt}{\baselineskip}\textbf{LISTA DE ILUSTRAÇÕES}}
\pdfbookmark[0]{\listfigurename}{lof}
\listoffigures*

\cleardoublepage

\tableofcontents*

\cleardoublepage
 
\mainmatter

%% Inclusão dos capítulos ao documento principal
\chapter{JUSTIFICATIVA}
\label{justificativa}

O meu pedido de prorrogação de bolsa de doutorado foi embasado nos artigos 2° e 4° da portaria 55
da COORDENAÇÃO DE APERFEIÇOAMENTO DE PESSOAL DE NÍVEL SUPERIOR - CAPES de 29 de Abril de 2020.
Esta portaria dispõe sobre a prorrogação excepcional
dos prazos de vigênciade das bolsas de mestrado e doutorado no país da CAPES
e estabelece os critérios necessários para o pedido de prorrogação do período de vigência das bolsas.
Em virtude da aceitação do meu pedido de prorrogação de bolsa de doutorado por parte da CAPES,
considero pertinente a utilização da mesma justificativa para o pedido de extensão do prazo da defesa da tese
por mais um ano além do período normal de defesa.

Como explicitado no Artigo 2° da Portaria 55 da CAPES abaixo:

"Art. 2º Fica autorizada, nos termos desta Portaria, a prorrogação dos prazos de vigência das bolsas de estudo de mestrado e doutorado concedidas no âmbito dos programas e acordos de competência da Diretoria de Programas e Bolsas no País da CAPES quando as restrições decorrentes do isolamento social necessário ao combate à pandemia da CoViD-19 tenham afetado o regular desenvolvimento do curso de pós-graduação ou o adequado desempenho dos mestrandos e doutorandos."

E explicitado no Art. 4°:

"Art. 4º São circunstâncias aptas a dar ensejo à prorrogação autorizada por esta Portaria:
I - o cancelamento ou o adiamento de atividades presenciais necessárias ao desenvolvimento do curso, que não possam ser supridas adequadamente por meio de ensino à distância ou outros meios, tais como atividades laboratoriais ou de campo, coleta de dados, entre outras; II - restrições temporárias de acesso a instalações necessárias ao desenvolvimento das atividades do curso;
III - outras situações que tenham imposto dificuldades não antevistas aos mestrandos e doutorandos, respeitados os limites fixados por esta Portaria."

O disposto nos artigos acima embasa a minha justificativa para o pedido de extensão do prazo de defesa da minha
tese de doutorado, pois durante o período de isolamento social me encontrei impossibilitado de desenvolver as atividades laboratoriais de que depende o desenvolvimento regular da minha tese de Doutorado. Haja vista que os experimentos computacionais que são a base dos resultados a serem apresentados na minha defesa não puderam ter a sua devida continuidade em virtude da drástica redução do poder computacional disponível para realizá-los. Cito alguns dos empecilhos abaixo:

\begin{itemize}


\item Os experimentos numéricos desenvolvidos no Laboratório de Inversão de Ondas Sísmicas (LIOS) dependem de estratégias de otimização e paralelização realizadas nas Virtual Central Processing Units (VCPUS) dos computadores do laboratório através do programa SCons e do pacote de processamento sísmico Madagascar.
Estes computadores possuem 12 VCPUS cada, 2 núcleos em cada um dos seus 6 processadores Xeon.

\item O Meu computador pessoal (notebook positivo stilo xci3650) dispõe de 2 núcleos do processador Intel Celeron. O ambiente do laboratório é refrigerado e utiliza a energia elétrica da UFPA. Um ambiente refrigerado é necessário para um maior desempenho dos experimentos numéricos realizados e estes podem rodar os algoritmos de inversão em looping por vários dias seguidos. O calor gerado reduz o desempenho da máquina e sua vida útil e o gasto com energia elétrica também deve ser considerado. Além disso, os computadores e o ambiente do laboratório são fabricados para este tipo de desenvolvimento de
alto desempenho e o meu notebook não é.

\item A atividade no laboratório permite realizar o desenvolvimento e a experimentação ao mesmo tempo. No LIOS eu realizava o desenvolvimento dos algoritmos e a escrita do código no meu notebook enquanto rodava os experimentos numéricos anteriormente citados nos computadores do laboratório, oque já não posso mais fazer.

\item Os computadores do LIOS permitem uma estratégia simulando a "clusterização". Estes computadores me permitiam rodar vários experimentos numéricos e testes de diferentes modelos sintéticos ao mesmo tempo, em diferentes máquinas, e com a estratégia de paralelização supracitada nas VCPUS de cada máquina (cada uma com 12 núcleos), controlando e balanceando a carga com o software de integração contínua "Jenkins" e acessando as máquinas via OpenSSH.
Oque não pode mais ser feito e custaria milhares de reais para reproduzir em algum serviço de VPS ou Cloud Computing.

\item Para que eu possa realizar as atividades citadas acima, necessárias ao desenvolvimento da minha Tese de doutorado, irei necessitar tanto do tempo extra quanto do recurso financeiro da minha bolsa de doutorado, como me dedico exclusivamente à pesquisa, esta representa a minha única fonte de renda.

\end{itemize}

A justificativa apresentada acima é suficiente e está contemplada pelos artigos supracitados da portaria 55. Esta justificativa é evidenciada pela lista de empecilhos apresentados. Estes contribuem para reduzir o desempenho do regular desenvolvimento da minha Tese de Doutorado pelos meses que durarem o isolamento social. Por conseguinte, como esta
justificativa foi considerada pertinente por parte da CAPES para a prorrogação do período de vigência da minha bolsa
de doutorado, considero também justificado o meu pedidido de prorrogação do prazo para a defesa da tese de doutorado por mais
um ano além do prazo normal de defesa.

% \chapter{INTRODUÇÃO}
\label{cap1}

A estratégia utilizada para a inversão do modelo de velocidades apresentada na
qualificação de soutorado \cite{relatorio} foi abandonada em prol de uma nova estratégia de
inversão baseada na NIP tomografia e na stereo tomografia \cite{niptomo,stereo}.
A inversão do modelo de velocidades a partir da simulação e focalização
de difrações sobre a seção empilhada \cite{relatorio}  
possui uma séria limitação: Esta necessita do modelo de velocidades
de sobretempo normal (informação a priori) para simular as difrações sobre a seção empilhada e realizar a inversão através da focalização destas difrações. O empilhamento Elemento de Reflaxão Comum (ERC)
não produz o modelo de velocidades de sobretempo normal,
a seção empilhada é obtida a partir dos parâmetros $R_{NIP}$ e $\beta_0$ do método 
Superfície de Reflexão Comum (SRC) de afastamento nulo.

Assim, neste relatório, propomos uma nova estratégia para a inversão do modelo de velocidades em
profundidade utilizando um método baseado na NIP tomografia e na stereo tomografia: A partir da seção
empilhada ERC obtida \cite{relatorio} realizamos o picking dos tempos de trânsito sobre os eventos de
reflexão na seção. Estes tempos de trânsito são utilizados para localizar fontes pontuais sobre os relfetores
em um modelo de velocidades de background. Isto é feito através do traçamento de raios normais a partir
da superfície de aquisição em direção ao modelo em profundidade. A direção inicial do raio normal
é dada pelo ângulo de incidência $\beta_0$.
Realizamos a modelagem direta traçando raios de reflexão a partir destas fontes pontuais em profundidade
até atingirem a superfície de reflexão e comparamos os tempos de trânsito obtidos com os tempos de trânsito
calculados através da fórmula do ERC. O pressuposto da nossa estratégia de inversão é que os tempos
de trânsito obtidos com o traçamento de raios serão próximos aos tempos de trânsito calculados para
o modelo de velocidades ótimo.





%\include{cre} % Método CRE fundamentos teóricos
% \include{vfsa} % Otimização global VFSA
% \include{pef} % Teoria intepolação com filtros PEF
% \include{modeling} % Resultado da modelagem
% \include{interpolacao} % Resultados da Interpoção com PEF
% \include{empilhamento} % Resultados do empilhamento cre
% \chapter{INVERSÃO DO MODELO DE VELOCIDADES}
\label{cap8:velocidades}

Propomos a seguinte metodologia para a inversão do modelo de velocidades para o cumprimento do objetivo da tese:
Utilizar a simulação de resposta de difrações \cite{diffractions} e a análise de velocidade pós empilhamento
com continuação de velocidades para obter as velocidades de migração que melhor focalizam as hipérboles
de difração da seção simulada \cite{sep_dif}. A otimização do modelo de velocidades se dará através de
um algoritmo iterativo com atualização do modelo de velocidades através do VFSA \cite{mesquita}.

A determinação da velocidade através da simulação de resposta de difração
consiste dos seguintes passos \cite{diffractions}:

\begin{enumerate}
 \item Determinação da velocidade NMO: Isto pode ser feito em cascata nos pontos da
seção empilhada, camada por camada otimizando a velocidade uma camada por vez.
\item Determinação das inclinações das curvas de reflexão: Na seção de afastamento nulo e no cubo de dados empilhados
Este passo será feito através da utilização de filtros de destruição de ondas planas \cite{planas}.
\item Transformação dos traços zero offset selecionados em respostas de pontos difratores em pares $m_0, t_0$
selecionados.
\item Aplicação da migração em profundidade pós empilhamento com análise de velocidades
em looping.
\end{enumerate}

A partir da seção empilhada ERC transformamos cada par $m_0, t_0$ sobre uma interface refletora em uma resposta simulada de
uma fonte pontual (Ponto difrator) na forma de hipérboles de difração. Escolhido um par $m_0, t_0$, dado um modelo inicial de
velocidades, simulamos uma resposta de difração espalhando a amplitude $A(m_0,t_0)$ na curva de 
tempo de trânsito de difração $T_D$ \cite{diffractions}:

\begin{equation}
 \label{eq:8.1}
 T_D^2 = [t_0 - 2 m (p_0 \cos \beta_0) ]^2 + \frac{4m^2}{v_{NMO}^2}
\end{equation}

Onde $t_0$ é o tempo do raio normal, $m$ é a coordenada do PMC, $p_0$ é a inclinação da curva de tempo de reflexão e $v_{NMO}$
é a velocidade de sobretempo normal na direção do PMC. A inversão do modelo de velocidades
é a busca pelas velocidades de migração que melhor focalizam as hipérboles de difração simuladas.
As Figuras \ref{fig:8.1}-\ref{fig:8.3} são um exemplo da metodologia de resposta de difrações simuladas na inversão do
modelo de velocidades de um modelo geológico sintético de uma estrutura salina.

\begin{figure}[htb]
\caption{Respostas de difração simuladas na seção empilhada em um modelo geológico sintético de uma estrutura salina.}
\begin{center}
\includegraphics[scale=0.30]{images/simula_diff.png}
\vspace{-0.3cm}
\end{center}
\begin{center}
 Fonte: \cite{diffractions}.
\end{center}
\label{fig:8.1}
\end{figure}

\begin{figure}[htb]
\caption{Passo intermediário na determinação da velocidade de migração em loop.}
\begin{center}
\includegraphics[scale=0.30]{images/simula_diff_campo.png}
\vspace{-0.3cm}
\end{center}
\begin{center}
 Fonte: \cite{diffractions}.
\end{center}
\label{fig:8.2}
\end{figure}

\begin{figure}[htb]
\caption{Imagem em profundidade obtido através da inversão a partir de respostas de difração simuladas 
em um modelo geológico sintético de uma estrutura salina.}
\begin{center}
\includegraphics[scale=0.30]{images/simula_diff_vel.png}
\vspace{-0.3cm}
\end{center}
\begin{center}
 Fonte: \cite{diffractions}.
\end{center}
\label{fig:8.3}
\end{figure}

A análise de velocidades automatizada no domínio pós empilhado e o imageamento de alta resolução
para heterogeneidades de pequena escala serão feitos através da produção de seções migradas no tempo 
com as hipérboles de difração focalizadas. 
Este processo contará com a análise automatizada de focalização para a detecção das
velocidades de migração ótimas no imageamento das hipérboles de difração simuladas 
utilizando a medida de variação máxima local \cite{sep_dif}:

\begin{equation}
 \label{eq:8.2}
 \phi_i = \frac{ N \sum^N_{i=1} s_i^4 }{ \left( \sum^N_{i=1} s_i^2 \right) } 
\end{equation}

Onde $s_i$ é a amplitude da amostra do sinal sísmico em uma janela de tamanho $N$.
Sinais bem focalizados possuem alta variação máxima local $\phi_i$. Desta forma, produziremos
os painéis de focalização de imagem calculando $\phi_i$ para cada ponto na seção migrada
e realizaremos a análise de velocidades automatizada (Um exemplo dos painéis de focalização de imagem
está na Figura \ref{fig:8.4}).

\begin{figure}[htb]
\caption{Exemplo de painéis de focalização de imagem. As cores representam os valores de variação
máxima local. O eixo das abcissas é a velocidade em Km/s. O eixo das ordenadas é o tempo em s.}
\begin{center}
\includegraphics[scale=0.30]{images/panelfig.png}
\vspace{-0.3cm}
\end{center}
\begin{center}
 Fonte: \cite{sep_dif}.
\end{center}
\label{fig:8.4}
\end{figure}

A principal diferença desta metodologia para a análise
de velocidades convencional é que a informação sobre a velocidade
será obtida da focalização das difrações ao invés dos painéis de coerência \cite{sep_dif}.

A migração no tempo será realizada a partir de um algoritmo de continuação de velocidades \cite{fomel2003a},
um método que performa
a análise de velocidade de migração em tempo continuando imagens sísmicas na velocidade, também chamado de
``ondas imagem'' \cite{hubral1996}. A teoria da continuação de velocidades mostra ser possível realizar a migração no tempo
com um conjunto de diferentes velocidades fazendo pequenos incrementos na velocidade diretamente no domínio da imagem.
Implementaremos através de um algoritmo baseado na tranformada rápida de Fourier \cite{bfomel2003}.

A Figura \ref{fig:8.5} apresenta o algoritmo proposto para a obtenção do modelo de velocidades. O algoritmo funciona em um
laço seguindo os passos anteriormente descritos e a atualização do modelo de velocidades da iteração é realizada com
auxílio do VFSA até atingir um modelo ótimo. O critério de convergência será a coerência entre os dados modelados a partir do
o modelo de velocidades da iteração atual e os dados sísmicos pré empilhados \cite{mesquita}.

\begin{figure}[htb]
\caption{Fluxograma do algoritmo de inversão do modelo de velocidades.}
\begin{center}
\includegraphics[scale=0.5]{images/fluxoVel.png}
\vspace{-0.3cm}
\end{center}
\begin{center}
 Fonte: Do Autor.
\end{center}
\label{fig:8.5}
\end{figure} % Proposta de uma metodologia para inversão do modelo de velocidades
% \include{cronograma} % cronograma de atividades e entrega do trabalho
% \include{conclusao} % conclusão discutindo os resultados obtidos e esperados

\bookmarksetup{startatroot}

\bibliography{mybib}

%% Inicia os apêndices
%% ---
% \begin{apendicesenv}

%% Imprime uma página indicando o início dos apêndices
% \partapendices
%\newpage\null\thispagestyle{empty}
% \begin{center}
% \vspace{10cm}
% {\large\textbf{APÊNDICES}}
% \end{center}
% \newpage

% \input{apendice_tempo_transito_cre}
% \input{apendiceB}
% \input{apendiceC}
% \input{apendiceD}
% \input{MT}
%
% \end{apendicesenv}

%% Inicia os anexos
%\begin{anexosenv}
%\partanexos
%\include{AnexoA}

%\end{anexosenv}

\cleardoublepage
\phantomsection 
\printindex

\end{document}
